\documentclass[oneside, 12pt]{book}
\usepackage[polish]{babel}
\usepackage{geometry}
\usepackage{lipsum}
\usepackage{nameref}
\usepackage{polski}
\usepackage[utf8]{inputenc}
\usepackage{setspace}
\usepackage[export]{adjustbox}
\usepackage{biblatex}
\usepackage{titlesec}
\usepackage{csquotes}
\usepackage{float}

\usepackage[table]{xcolor} % For cell colors
\usepackage{array} % For table alignment
\usepackage{colortbl} % For color in tables
\usepackage{booktabs} % For better-looking tables
\usepackage{svg}

\usepackage[pdfsubject={Pielat K.: Analysis of the Influence of Acoustic
Parameters and Lyrics on the Characteristics of Music Tracks Using AI Methods},
pdfauthor={Pielat K.},
pdftitle={Pielat K.: Analysis of the Influence of Acoustic Parameters and
Lyrics on the Characteristics of Music Tracks Using AI Methods},
pdfkeywords={musical analysis; nlp; acoustic analysis}]{hyperref}
\usepackage{tabularray}

\addto\captionspolish{
  \renewcommand{\figurename}{Figure}
}

\addto\captionspolish{
  \renewcommand{\tablename}{Table}
}

\addto\captionspolish{
  \renewcommand{\contentsname}{Table of Contents}
}

\setcounter{tocdepth}{3}
\setcounter{secnumdepth}{3}

\addbibresource{bibliography.bib}

\titleformat{\chapter}[block]
  {\normalfont\huge\bfseries}{\thechapter.}{1em}{\Huge}

\titlespacing*{\chapter}{0pt}{50pt}{40pt}

\geometry{
    top=1in,
    bottom=1in,
    outer=1in,
    inner=1in,
}

\pagestyle{plain} %Numeracja na dole

\linespread{1.5} % Interlinia 1.5

\setlength{\parindent}{1.25cm} % Wcięcie akapitu 1.25cm

\begin{document}
\thispagestyle{empty}

% Informacje na stronę tytułową
\newcommand{\autor}{\fontfamily{qhv}\fontshape{n}\selectfont\LARGE\bfseries
Krystian Pielat}
%
\newcommand{\album}{\fontfamily{qhv}\fontshape{n}\selectfont
143007}

\newcommand{\pltitle}{\fontfamily{qhv}\fontshape{n}\selectfont\Large\bfseries
Analiza wpływu parametrów akustycznych i tekstów na charakterystykę utworów muzycznych z wykorzystaniem metod sztucznej inteligencji}
\newcommand{\engtitle}{\fontfamily{qhv}\fontshape{n}\selectfont\Large\bfseries
 Analysis of the Influence of Acoustic Parameters and Lyrics on the Characteristics of Music Tracks Using AI Methods}
\newcommand{\level}{\fontfamily{qhv}\fontshape{n}\selectfont\Large
inżynierska} %jeśli magisterka to trzeba sobie przystosować pamiętając o specjalności
\newcommand{\kierunek}{\fontfamily{qhv}\fontshape{n}\selectfont\Large
Informatyka}
%
\newcommand{\promotor}{\fontfamily{qhv}\fontshape{n}\selectfont\large\bfseries
dr inż. Daniela Grzonki}
\newcommand{\spec}{\fontfamily{qhv}\fontshape{n}\selectfont\Large\bfseries
Brak}
\newcommand{\rok}{\fontfamily{qhv}\fontshape{n}\selectfont
2024}

% Strona tytułowa
\begin{titlepage}
\thispagestyle{empty}
\vspace*{1ex}
\fontfamily{qhv}\fontshape{n}\selectfont
\hspace{-3.5em}\parbox{0.15\textwidth}{\includegraphics[height=0.15\textwidth]{logo_pk/logoPK.png}}
\parbox[c]{0.7\textwidth}
{\centering
\Large
{\bfseries Politechnika Krakowska}\\
{\bfseries im. Tadeusza Kościuszki}\\

\smallskip

\Large Wydział Informatyki i Telekomunikacji}\hspace{0.3em}
 \parbox{0.2\textwidth}{\includegraphics[width=0.15\textwidth]{logo_pk/logoIT.png}}%\hspace{0.12em}
\parbox[c]{0.15\textwidth}



 \vspace{0.10\textheight}

 \center{\autor}
 \smallskip
 \center{\large numer albumu: \album}

\vspace*{2ex}

 \center{\pltitle}

 \smallskip

 \center{\engtitle}

 \bigskip

 \center{\Large\bfseries Praca \level \\[0.2ex] %jeśli magisterka to trzeba sobie przystosować pamiętając o specjalności
\fontfamily{qhv}\fontshape{n}\selectfont\Large\bfseries na kierunku
 \kierunek\\[0.2ex]
 }

\vspace{0.11\textheight}

\hspace{0.57\textwidth} \parbox{0.42\textwidth}{Praca przygotowana pod
kierunkiem:
\\
%
\bfseries \promotor}

\vfill

 % wpisać rok
 \center{Kraków, \rok}
 \end{titlepage}

 \newpage

% \vspace*{17cm}
% \begin{flushright}
% \textit{Tu można sobie coś wpisać}
% \end{flushright}






\chapter*{Abstrakt}
Niniejsza praca bada zależności pomiędzy metadanymi, tekstami i parametrami
akustycznymi utworów muzycznych wykorzystując zaawansowaną analizę danych i
metody sztucznej inteligencji. Opracowano zróżnicowany zbiór danych utworów
składający się z danych platformy streamingowej Spotify, właściwości
akustycznych, tekstów i plików mp3, wykorzystując dobrze zdefiniowaną
metodologię zbierania danych, aby analizować wzorce i interakcje związane z
popularnością, sentymentem, gatunkiem i innymi właściwościami utworów. Badania
uwzględniają również, jak te cechy zmieniają się na przestrzeni dekad.
Wykorzystując dogłębną eksploracyjną analizę danych, rozbudowana inżynierię
cech, techniki przetwarzania języka naturalnego i uczenia maszynowego,
modelowanie tematów oraz metody statystyczne praca ta ma na celu odkrycie
ukrytych wzorców, weryfikację hipotez oraz zrozumienie interakcji obecnych w
zebranych danych.

% Niniejsza praca bada zależności pomiędzy parametrami akustycznymi, tekstami i 
% cechami utworów muzycznych przy użyciu metod sztucznej inteligencji. Opracowano
% i przeanalizowano zbiór danych utworów, obejmujący metadane Spotify, cechy
% akustyczne, teksty piosenek oraz ścieżki mp3, w celu zbadania wzorców
% związanych z tożsamością gatunkową, sentymentem i innymi właściwościami
% utworów. Badania uwzględniają również, jak te cechy zmieniają się w czasie.
% Wykorzystując rozbudowaną inżynierię cech, dogłębną eksploracyjną analizę
% danych, przetwarzanie języka naturalnego, techniki uczenia maszynowego,
% modelowanie tematów oraz metody statystyczne, praca analizuje relacje między
% elementami muzycznymi i tekstowymi. Opracowano i zinterpretowano modele
% predykcyjne za pomocą metod wyjaśnialnej sztucznej inteligencji, co zapewnia
% przejrzystość i dostarcza cennych wniosków na temat branży muzycznej.

\chapter*{Abstract}
This thesis explores the relationships between metadata, lyrics and acoustic
characteristics of music tracks using advanced data analysis and various
artificial intelligence methods. It creates a diverse dataset of songs
consisting of Spotify data, audio features, lyrics and mp3 tracks using a
well-defined data collection methodology to study patterns and interactions
related to song popularity, sentiment, genre identity and other song
characteristics. This study also examines how these properties evolve over
time. Using thorough exploratory data analysis, comprehensive feature
engineering, natural language processing and machine learning techniques, topic
modeling and statistical methods it aims to uncover hidden patterns, verify
hypotheses and understand interactions present in the collected data.


% This thesis explores the relationships between acoustic parameters, lyrics, and
% the characteristics of music tracks using artificial intelligence methods. It
% creates and analyzes a dataset of songs, including Spotify metadata, audio
% features, lyrics, and mp3 tracks, to study patterns related to song popularity,
% genre identity, sentiment, and other song characteristics. The research also
% examines how these properties evolve over time. Using comprehensive feature
% engineering, thorough exploratory data analysis, natural language processing,
% machine learning techniques, topic modeling, and statistical methods, it
% investigates connections between musical and lyrical elements. Predictive
% models are developed and interpreted with explainable AI methods to ensure
% transparency, providing valuable insights into the music industry.
%





\tableofcontents
\clearpage

\chapter{Introduction}
\label{cha:introduction}

\section{Problem Statement}
\label{sec:problemstatement}
Music has been a prominent part of human culture for generations. It's been
evolving alongside societies, reflecting their creativity, emotions, values and
emerging trends. In this day and age music has become more accessible and
popular than ever, especially thanks to advancements in technology. Rapidly
expanding global audience, growing number of artists and overwhelming number of
songs released every day presents us with a great opportunity to investigate
more closely its characteristics through the lens of data-driven analyses and
machine learning techniques. 

Based on the acoustic features that can be extracted from the music tracks, and
textual features extracted from song lyrics and information available on
Spotify, this paper attempts to understand better complex relationships between
different music track characteristics and uncover insights into how music is
perceived by listeners, the factors that influence song's popularity, the
defining traits of various music genres and how trends shape the evolution of
music and its characteristics over the years. 

This approach offers a modern, quantitative perspective on music, enabling us
to further explore relationships between different musical and textual
characteristics.




% \begin{center}
% \begin{figure}[ht]
%   \centering
%   \includegraphics[width=6in]{img/plik.png}
%   \caption{Podpis pod rysunkiem\cite{nature_behavior}.}
%   \label{Figure:fig_beh}
% \end{figure}
% \end{center}

%---------------------------------------------------------------------------

\section{Research Questions \& Objectives}
\label{sec:researchquestions}
% This thesis aims to leverage a diverse dataset of songs collected through a
% systematic metehodology, consisting of metadata, Spotify's audio features,
% lyrics and mp3 tracks. 

The aim of this thesis is to build a diverse dataset of songs with lyrical,
acoustic and metadata features through a well-defined data collection
methodology and use various advanced data analysis methods and explainable
artificial intelligence(\textit{XAI}) to:

\begin{enumerate}
  \item Explore the relationships between lyrical, acoustic, and metadata
    features of songs to uncover meaningful patterns and dependencies.
  \item Analyze defining characteristics of songs across different genres and
    investigate how these features contribute to genre identity and what are
    defining traits for each genre included in the dataset.
  \item Develop predictive models to estimate various song attributes (e.g.
    popularity, explicitness, sentiment) and use XAI methodologies to interpret
    model's predictions and extract meaningful insights and relationships.
  \item Identify commonly recurring themes in lyrics and explore their
    significance and distinctive traits using topic modelling. 
  \item Uncover temporal trends and shifts in music characteristics over time.
\end{enumerate}

Additionally, this thesis aims to create a practical and reusable explainable
AI(\textit{XAI}) evaluation framework and a clear and systematic method for
collecting data. These tools are designed to make it easier to build on this
work, apply the methods to new projects and research.


%---------------------------------------------------------------------------


\section{Thesis Structure}
\label{sec:thesisstructure}

This thesis begins with an introduction and explaination of research objectives.
It's followed by a discussion of related studies, summarizing each of them and
explaining how  they relate to this work. 

The data chapter explains the sampling method that was used to select the songs
from Spotify, the data collection framework, the sources for the data and the
data itself.

Next, the methodology chapter outlines the feature engineering process,
describing all extracted features, and introduces all the techniques
and methods used in this study. 

The chapter that comes after that describes the extensive exploratory data
analysis conducted on the collected data with the aim of exploring relationships
between features and identifying distinctive traits of each genre.

The experiments chapter presents all the experiments conducted in this study
and discusses their results. It also compares them to different benchmarks and
attempts to interpret the results.

The thesis concludes with a summary of contributions, a description
of the limitations, key findings, and recommendations for future work. 

\chapter{Related Work}
\label{cha:literaturereview}

%---------------------------------------------------------------------------

\section{Acoustic Features in Music Analysis}
\label{sec:acousticfeaturesinmusicanalysis}

\subsubsection*{``Music Genre Classification Using MFCC, k-NN, and SVM
Classifier''}


This study explores music genre classification using MFCCs and Chroma features
on the GTZAN dataset, that consists of 900 tracks across 9 genres. The
best-performing model was an SVM with a polynomial kernel, achieving accuracy
of 78\%. Some genres were identified to have overlapping characteristics which
posed problems for the classification model. \cite{music_genre_classification_mfcc}

The paper highlights the effectiveness of MFCCs and Chroma features for
audio-based classification, aligning with this thesis's use of acoustic
features. However, unlike this study, the thesis extends the analysis by
incorporating lyrical features, metadata, and audio features provided by
Spotify, addressing limitations in feature diversity.

\subsubsection*{``Classifying Music Audio with Timbral and Chroma Features''}

The paper ``Classifying Music Audio with Timbral and Chroma Features'' examines
the use of timbral (e.g., MFCCs) and harmonic (e.g., chroma) features in music
classification tasks. The study highlights that MFCCs, which represent timbral
qualities, are effective for tasks like \textbf{artist identification, achieving 56\%
accuracy in a 20-way classification task}. When combined with beat-synchronous
chroma features, which capture harmonic and melodic information, the model's
accuracy improved to 59\%. This demonstrates the importance of leveraging
complementary audio features for more robust classification. These findings
emphasize the potential of combining diverse acoustic features in predictive
models, aligning closely with the objectives of this thesis. \cite{classifying_music_audio}



%---------------------------------------------------------------------------

\section{Lyrical Features in Music Analysis}
\label{sec:lyricalfeaturesinmusicanalysis}

\subsubsection*{``Using Machine Learning Analysis to Interpret the Relationship
Between Music Emotion and Lyric Features''}


This study investigates the relationship between lyrical features and perceived
music emotions using 2,372 Chinese pop songs. Lyric features were extracted
with LIWC (Linguistic Inquiry and Word Count), and audio features
such as MFCCs and chroma were derived using Librosa. The analysis highlighted
that lyrical features like the frequency of positive and negative emotion words
contributed significantly to predicting the perceived valence of music, whereas
audio features were dominant in predicting perceived arousal. \cite{valence_and_lyrics}

The study utilized Random Forest regression models, demonstrating that
combining lyric and audio features improved \textbf{predictions of valence,
with an $R^2$ of 0.481}, but lyrics had little impact on arousal models. These
findings align with this thesis's use of both lyrical and acoustic features for
predictive tasks but differ in the application of explainable AI (XAI)
techniques like SHAP for feature interpretation. Moreover, this thesis expands
this study by incorporating Spotify audio features and metadata for broader
analytical capabilities and more diverse research objectives.


\subsubsection*{``Sentiment Analysis and Lyrics Theme Recognition Using NLP Techniques''}


This paper investigates the relationship between sentiment and themes in music
lyrics using Natural Language Processing (NLP) techniques. It applies sentiment
analysis and thematic recognition across a diverse dataset, identifying
emotional nuances and recurring topics in the lyrics. The analysis identifies
correlations between sentiment categories (positive, negative, neutral) and
thematic clusters (e.g., love, social justice, personal reflection). The
authors use techniques like Latent Dirichlet Allocation (LDA) for topic
modeling and Support Vector Machines (SVM) for sentiment classification. \cite{du_2024}

In the context of this thesis, this study aligns closely with the focus on
lyrical analysis, particularly in employing NLP-driven sentiment and thematic
classification. However, while this work emphasizes standalone lyric-based
analysis, this thesis extends the methodology by integrating Spotify’s audio
features, acoustic analysis, and metadata. 

%---------------------------------------------------------------------------

\section{Machine Learning in Music Analysis}
\label{sec:machinelearningfeaturesinmusicanalysis}

\subsubsection*{``Beyond Beats: A Recipe to Song Popularity? A Machine Learning Approach''}

The paper ``Beyond Beats'' investigates the predictive power of various machine
learning models for song popularity on a dataset of 30,000 songs spanning six
genres. It focuses on metadata and audio features fetched from
Spotify (\textit{danceability}, \textit{acousticness} etc.), as well as the
genre. The best performing model was \textbf{Random Forest and achieved 16.31 MAE.}
Predictions across all methods remained relatively modest, reflecting the
complex and multi-dimensional  nature of song popularity. Those results can be
greatly improved using additional features. The authors noted that predictive
accuracy was constrained by the absence of post-release factors such as
marketing, social media reception, and artist reputation. \cite{beyond_beats}


\subsubsection*{``Predicting Song Popularity in the Digital Age Through
Spotify’s Data''}  
Similarily to the previous one, the paper ``Predicting Song Popularity in the
Digital Age Through Spotify’s Data'' explores the relationship between
Spotify's audio features  and song popularity, using a dataset spanning from
1986 to 2022. The study employed linear regression to predict popularity and
\textbf{achieved an adjusted $R^2$ of 0.38}, highlighting the moderate
predictive power of these features. The analysis revealed that attributes such
as \textit{danceability} and \textit{duration} positively correlate with
popularity, while \textit{speechiness} tends to have a negative
impact. \cite{predicting_song_popularity_2024}

\chapter{Data}
\label{cha:data}
%---------------------------------------------------------------------------

\section{Dataset Description}
\label{sec:datasetdescription}

The dataset consists of around 6000 songs. For each song metadata, audio
recording, lyrics and spotify audio features were fetched.

\section{Data Collection Methods}
\label{sec:datacollectionmethods}

\subsection{Sources}
The data was collected from various sources. Metadata and audio features
were fetched from Spotify API. The lyrics were scraped from MusixMatch,
LetrasMus, MakeItPersonal or Lyrics Fandom, or fetched via the Genius API,
depending on  the availability. Audio files were downloaded from YouTube and
saved as mp3 files.

\subsection{Methodology}
The data collection process was automated for efficiency. The starting point
was a list of Spotify playlist URIs. Each playlist was processed sequentially. 
The script fetched metadata and audio features for each song in the
playlist.

Lyrics were then searched based on \textit{artist name} and \textit{title} of
each song, using multiple lyrics providers. The system continued to query
different sources until it found lyrics in at least one of them. If it failed
to retrieve lyrics for the song, it was discarded and wouldn't make it to the
final dataset.

Finally, the script  searched YouTube for each song and downloaded the first
relevant result, saving it to an mp3 file. All data was saved into a CSV
file named after the playlist and stored alongside the recordings.

The entire process was parallelized to significantly enhance the speed of data
acquisition. It was designed in a robust manner, with careful error handling
and ability to stop the process at any  time and pick up where it left off.

Additionally, to optimize performance, the entire process was parallelized,
significantly increasing the speed of data acquisition. It was designed with
robust error handling, ensuring data correctness and completeness. Moreover, it
featured the ability to pause and seamlessly resume the process, continuing
exactly where it left off without data loss.


%---------------------------------------------------------------------------


\subsection{Metadata}
\label{sec:metadata}
Song's information downloaded from Spotify API. The information it includes is:

\begin{itemize}
  \item \textbf{Popularity} - relative measure with values ranging from 0 to
    100 describing how popular the song is, estimated mostly based on total
    number of plays and how recent  those plays are
  \item \textbf{Explicitness} - whether or not the song contains explicit lyrics
  \item \textbf{Genre} - main genre of the artist(Spotify does not provide
    information about genre of each musical track)
  \item \textbf{Album Release Year} - the release year of the album that the
    song originates from
\end{itemize}

\textit{popularity}, \textit{explicitness},
\textit{genre}, and \textit{release year}. Since Spotify does not provide
information about the genre of specific songs, the genre extracted is the main
genre of the primary artist in the song.


\subsection{Spotify Audio Features}
\label{sec:spotifyaudiofeatures}
Those features were fetched from the Spotify API.They describe different
acoustic properties songs:
\begin{itemize}
  \item \textbf{Speechiness} - relative measure of spoken words in a track
  \item \textbf{Acousticness} - a confidence measure of whether the song is
    acoustic
  \item \textbf{Danceability} - a measure of how suitable the song is for
    dancing. Its based on parameters like tempo, rhythm stability, beat
    strength etc.
  \item \textbf{Energy} - a measure of perceived intensity of songs. Energetic
    tracks are usually louder, faster, feel more intense.
  \item \textbf{Loudness} - overall loudness of the  track in dB averaged
    across the entire track
  \item \textbf{Valence} - relative measure describing musical positiveness of
    a track
  \item \textbf{Instrumentalness} - measure of how likelihood of the track not
    containing vocals. In this paper since lyrics are mandatory it's used to
    discard instrumental tracks.
  \item \textbf{Liveness} - probability of the song being recorded during a
    live performance.
  \item \textbf{Key} - the key of the song, e.g. C\#.
  \item \textbf{Mode} - indicates the modality of a track(major / minor)
  \item \textbf{Tempo} - estimated tempo of a track in beats per minute(BPM)
  \item \textbf{Time Signature} - specifies how many beats there are in each
    bar
  \item \textbf{Duration} - the duration of the track in milliseconds
\end{itemize}


\subsection{Lyrics Features}
 The lyrics serve as a textual representation of the song's thematic,
 emotional, and linguistic elements. Since they came from various data sources,
 they had  to undergo cleaning procedure in order to remove faulty information
 and prepare them to be processed by the textual feature extraction class. The
 process consisted of:
 \begin{itemize}
  \item \textbf{Standardization} - lyrics were converted to lowercase
  \item \textbf{Noise Removal} - unnecessary characters, numbers and
    punctuation, as well as additional comments used by lyrics providers(e.g.
    'chorus') were removed
  \item \textbf{Stopwords Filtration} - exclusion of frequently occuring words
    that carry little information, like 'the' in English
  \item \textbf{Stemming} - words were reduced to their root forms to enhance
    uniformity and reduce corpus size
 \end{itemize}

 This process laid foundation for further extraction of textual features used
 for exploratory data analysis, statistical inference and training ML models.
 That process will be explained in later chapter.


%---------------------------------------------------------------------------

\section{Tools and Libraries Used}
\label{sec:toolsandlibrariesused}
Python libraries used to facilitate the data acquisition were:
\begin{itemize}
  \item \textit{Spotipy} - a lightweight python library for Spotify API
  \item \textit{youtube-dl} - a library used to find and download youtube videos
  \item \textit{BeautifulSoup} - a library used for extracting information from
    HTML, commonly used for web scraping
\end{itemize}




\chapter{Methodologies}
\label{cha:methodologies}


This chapter explains the methodologies used to address the research objectives
in this thesis. It describes the feature extraction and  engineering methods,
the use of explainable AI methods, model training and optimization approaches,
clustering and dimensionality reduction techniques and statistical hypothesis
testing methods.
%---------------------------------------------------------------------------

\section{Feature Engineering}
\label{sec:featureengineering}

Feature engineering is a critical step in the research process, as it involves
transforming raw data acquired using the data collection script into meaningful
representations that can be analyzed or used for predictive modeling. This
section describes the methodologies used to extract acoustic features from the
mp3 files and lyrical features from the lyrics. These features are designed to
capture key characteristics of the songs, enabling deeper insights into their
patterns and relationships.

\subsection{Acoustic Features}
\label{sec:acousticfeatures}

 These features provide a quantitative representation of the audio properties
 of each song and were extracted directly from the audio files in MP3 format.
 They describe various aspects of audio signal and provide insights into the
 rhythm, timbre, harmony and other acoustic properties. The extraction was done
 using \textit{Librosa} and was automated and parallelized to make it suitable
 for processing large amounts of data. 

\subsubsection*{MFCC - Mel Frequency Cepstral Coefficients}
MFCCs represent the short-term power spectrum of a song on a mel-scale and are
widely used for timbre analysis. These coefficients capture the tonal quality
of the audio and help differentiate between different instruments and vocal
characteristics.


\subsubsection*{Chroma}
Chroma vectors represent the intensity of each pitch class (e.g., C, C\#, D,
etc.) in the audio. These features provide a harmonic representation of the
song and are useful for analyzing chord progressions and harmonic structures.

\subsubsection*{Spectral Contrast}
Spectral contrast measures the difference in amplitude between peaks and
valleys in the spectrum. It provides insights into the harmonic and timbral
content of a song, particularly useful for distinguishing between smooth and
complex textures.

\subsubsection*{Other Features}
Two additional features were extracted:
\begin{itemize}
  \item \textbf{Tempo} - refers to the speed of the song, measured in \textit{Beats Per Minute(BPM)}
  \item \textbf{Zero Crossing Rate(ZCR)} - measures the rate at which the audio
    signal changes sign. It's commonly used as a measure of noisiness or
    percussive nature of signal
\end{itemize}

%---------------------------------------------------------------------------

\subsection{Lyrical Features}
\label{sec:lyricalfeatures}

Lyrical features  were extracted from the lyrics fetched during the data
collection process. They aim to provide a linguistic and semantic
representation of the track, capturing their complexity, sentiment and
stylistic  attributes. The cleaning and extraction  process utilized various
NLP libraries like \textit{NLTK}, \textit{spaCy} and \textit{TextBlob},
alongside with custom ad-hoc algorithms. Similarily to acoustic features the
implementation allowed for simple and intuitive usage under clear and
comprehensible interface, with parallelization of the computation process for
increased performance. The features extracted can be grouped as follows:

\subsubsection*{Basic Linguistic Metrics}
\begin{itemize}
  \item \textbf{Unique Word Count} - measures number of unique words in the lyrics, indicating diversity
  \item \textbf{Type-Token Ratio} - a measure of lexical richness: ratio of unique words to total words
  \item \textbf{Word Count} - total  number of words, baseline for text size and compexity
  \item \textbf{Noun and Verb Ratios} - proporrtios of nous and vers relative to the total word count
\end{itemize}


\subsubsection*{Sentiment and Emotional Tone}
\begin{itemize}
  \item \textbf{Sentiment Polarity} - a measure of overall sentiment(positive
    vs. negative) of the text
  \item \textbf{Sentiment Subjectivity} - represents the degree of subjectivity
    in the lyrics, attempting to make a distinction between factual and
    opinionated content
  \item \textbf{VADER Compound} - a sentiment score derived from the VADER tool
  \item \textbf{Sentiment Variability} - standard deviation of sentiment on
    subsets of lyrics, a metric  aiming to capture fluctuations of sentiment
    throughout the song, highlighting emotional complexity
  \item \textbf{} -
\end{itemize}


\subsubsection*{Stylistic Features}
\begin{itemize}
  \item \textbf{Repetition Count} - the frequency of repeated words
  \item \textbf{Rhyme Density} - a measure of how often rhymes occur in the
    text
  % \item \textbf{Linguistic Uniqueness} - a measure of 
\end{itemize}


\subsubsection*{Semantic and Complexity Features}
\begin{itemize}
  \item \textbf{Semantic Depth} - represents the richness and variety of
    meaning conveyed by the lyrics
  \item \textbf{Syntactic Complexity} - captures the sophistication of
    sentence structures
  \item \textbf{Lexical Richness} - quantifies the variety and richness of the
    vocabulary
\end{itemize}



\subsubsection*{Readability and Accessibility}
\begin{itemize}
  \item \textbf{Flesch Reading Ease} - indicates how easy the lyrics are to
    read
  \item \textbf{Gunning Fog} - a metric that estimates the  years of education
    required to understand the text
  \item \textbf{Dale Chall} - a metric that accounts for familiar and
    unfamiliar words in the text
\end{itemize}


\subsubsection*{Contextual Information}
  In process of feature extraction the \textbf{language} of lyrics was also
  identified using \textit{langdetect} library that uses a classification model
  to make predictions based on n-grams extracted from the text. The identified
  language was also used in the cleaning process, to identify which stemmer and
  stopwords language  to use.


\subsubsection*{TF-IDF (Term Frequency - Inverse Document Frequency)}

TF-IDF\cite{tfidf} is a measure that can quantify the relevance of tokens in a document
amongst a collection of documents. In the context of this study, it determines
how important a word in a song's lyrics is  compared to all other song's lyrics
in the dataset. It's used to highlight words that are unique or meaningful
while giving less importance to very common words like 'the' or 'and'.

It can be broken down into two parts:


\begin{itemize}
  \item \textbf{TF - Term Frequency} - Measures the frequency of a term within
    a document:
  \[
  TF(w, d) = \frac{f_{w, d}}{N_d}
  \]
  where:
  \begin{itemize}
    \item \( f_{w, d} \): The number of times the word \( w \) appears in the document \( d \).
    \item \( N_d \): The total number of words in the document \( d \).
  \end{itemize}

  \item \textbf{IDF - Inverse Document Frequency} - Measures the rarity of a
    term across a collection of documents:
  \[
  IDF(w) = \log{\frac{N}{1 + n_w}}
  \]
  where:
  \begin{itemize}
    \item \( N \): The total number of documents in the corpus.
    \item \( n_w \): The number of documents containing the word \( w \).
  \end{itemize}
\end{itemize}

%---------------------------------------------------------------------------

\section{Explainable AI Methods}
\label{sec:explainableaimethods}


Explainable AI (XAI) techniques provide insights into the decision-making
processes of ML models, making it possible to understand the complex
relationships captured within the training data. By bridging the gap between
the pattern-recognition capabilities of these models and their practical
applications, XAI enables transparency and improves the interpretability of
results.

In this study this methodology was applied in various experiments to understand
how specific variables influence others, with aim of uncovering rrelationships
in the data and validating hypotheses. A deeper  understanding of factors
driving model predictions ensured that the results were both reliable and
meaningfully adressed the research objectives. The  techniques employed in this paper are: 

\subsubsection*{SHAP (SHapley Additive exPlanations)}
SHAP values offer a detailed  breakdown of how individual features contribute
to  each prediction. It uses game theory principles to compute the importance
of each feature for a given output, providing both global insights on general
feature importance and  local explainations, showing how diffrenet features
contributed to specific predictions. In this study SHAP values were used in
order to explain the predictive process of trained Catboost models, allowing
for identification of key features for specific prediction task and
visualization of relationships.

\begin{center}
\begin{figure}[ht]
  \centering
  \includegraphics[width=4in]{img/shap_intro.png}
  \caption{SHAP}
  \label{Figure:fig_beh}
\end{figure}
\end{center}

\begin{center}
\begin{figure}[ht]
  \centering
  \includegraphics[width=4.5in]{img/shap_beeswarm.png}
  \caption{Example SHAP beeswarm plot showing impact of some lyrical features
  on the classifier of \textit{explicitness}}
  \label{Figure:fig_beh}
\end{figure}
\end{center}

\begin{center}
\begin{figure}[ht]
  \centering
  \includegraphics[width=3in]{img/shap_feature_importance.png}
  \caption{Example SHAP feature importance plot showing impact of some lyrical features
  on the classifier of \textit{explicitness}}
  \label{Figure:fig_beh}
\end{figure}
\end{center}
%---------------------------------------------------------------------------

\subsubsection*{Machine Learning Models}
In this study, CatBoost, a widely recognized gradient boosting algorithm known
for its high predictive accuracy and efficiency, was employed to build
classification and regression models. These models utilized a combination of
acoustic, lyrical and metadata features in order to predict the variable of
interest, leveraging Catboost's strengths in handling diverse and complex
datasets. The result models were then subjected to SHAP analysis, in order to
understand their decision process and understand the interactions between
features and the target variable. Catboost was chosen for its flexibility, ease
of use, robust performance, compatibility with SHAP and built-in support for
categorical features,  which eliminated the need for extensive preprocessing.

In order to further optimize the performance of these models, the
hyperparameter tuning library \textit{Optuna} was used. Its efficient
optimization framework allowed for systematic exploration of different sets of
hyperparameter configurations, ensuring the model achieved optimal performance while avoiding
overfitting. 

To address the challenges commonly encountered when training ML models on
complex datasets, following techniques were employed:
\begin{itemize}
  \item \textbf{Cross-validation} - cross-validation was used to reduce the
    risk of overfitting and provide reliable performance metrics. By dividing
    the data into multiple folds, the model was iteratively trained and
    validated on different subsets, therefore ensuring robust evaluation across
    the dataset and improved model's reliability, at the cost of increased
    computational time.
  \item \textbf{Class Weights} - to handle class imbalance in target variable
    labels, CatBoost offers a built-in capability to assign different penalties
    for misclassifications of specific classes. This adjustment improves
    model's ability to make accurate predictions across all classes, instead of
    favouring the majority class.
  \item \textbf{Out of sample evaluation} - model's performance was assessed on
    a separate test dataset that was excluded from the training. This step
    provided a reliable measure of model's ability to generalize on unseen data
    and ensured evaluation metrics  reflected its real predictive performance.
\end{itemize}

\section{Dimensionality Reduction - Principal Component Analysis (PCA)}
\label{sec:dimensionalityreduction}

Principal Component Analysis (PCA) reduces the number of features in large
datasets by transforming them into principal components that retain most of the
original information. It achieves this by converting potentially correlated
variables into a smaller set of less correlated variables, called principal
components, in a way that preserves as much of the original variance as
possible\cite{pca}. PCA is often employed to reduce dataset dimensionality and improve
generalization by reducing noise and redundancy in the data.

In this study, PCA was applied to the TF-IDF vectors derived from song lyrics.
TF-IDF vectors are typically high-dimensional, with thousands of features
representing individual terms across the corpus. Such high-dimensional data can
pose challenges, including increased computational complexity and a higher risk
of overfitting in machine learning models.

The use of PCA on these vectors reduced their dimensionality while preserving
as many significant patterns from the original vectors as possible. This
process improved computational efficiency and the interpretability of the data,
which is particularly important in the context of Explainable AI (XAI)
methodologies.

\begin{center}
\begin{figure}[ht]
  \centering
  \includegraphics[width=4in]{img/pca.png}
  \caption{A scatterplot showing the relationship between PC1 and PC2 when PCA
  is applied to a dataset. PC1 and PC2 axis are perpendicular to each other.\cite{pca}}
  \label{Figure:fig_beh}
\end{figure}
\end{center}


\section{Topic Modelling - Latent Dirichlet Allocation (LDA)}
\label{sec:topicmodelling}

Latent Dirichlet Allocation (LDA)\cite{lda} is a generative probabilistic model
designed to uncover latent topics within a collection of discrete data, such as
text corpus.  It represents each document as a mixture of topics, where each
topic is characterized by a distribution over words. In this study, LDA was
applied to the lyrics dataset to identify prevalent themes and topics across
different songs. By representing each song as a distribution over topics, LDA
provided insights into the thematic content of the lyrics, allowing for the
analysis of how these themes correlate with acoustic features and other song
attributes.

In this study, LDA was applied on the lyrics in order to identify commonly
ocurrirng topics in songs. Each song's lyrics were represented as a combination
of topics, and the most representative words for each topic were extracted.
This provided insights into the thematic content of lyrics, enabling further
exploration of relationships between lyrical topics and other features, like
popularity orr acoustic properties.

By combining the topics derived with LDA with additional features, such as
acoustic parameters, the study aimed to analyze the interplay between a song's
musical and lyrical components.


\section{Statistical Methods}
\label{sec:statisticalmethods}

Statistical methods provided a foundation for analyzing complex relationships
between features and a framework for descriptive data analysis and hypothesis
testing.

\subsection{Person Correlation}

Person Corrrelation is a statistical measure used to quantify linear
relationship between two continuous variables. The resulting coefficient ranges
from -1 to 1 and is computed in the following way:

\[
r = \frac{\sum{(x_i - \bar{x})(y_i - \bar{y})}}{\sqrt{\sum{(x_i - \bar{x})^2} \sum{(y_i - \bar{y})^2}}}
\]

Where:
\begin{itemize}
    \item \( x_i \) and \( y_i \): The data points for the two variables.
    \item \( \bar{x} \) and \( \bar{y} \): The mean values of the variables.
\end{itemize}

Interpretation:
\begin{itemize}
  \item An \textit{r} value close to 1 indicates very strong positive linear relationship
  \item An \textit{r} value close to -1 indicates very strong negative linear relationship
  \item An \textit{r} value close to 0 indicates little to no relationship
\end{itemize}



\subsection{Bootstrap Testing}

Bootstrap testing was used to verify hypotheses specified in the research
objectives. It's a resampling-based statistical technique that estimates the
variability of a statistic(e.g. mean or median)  by
repeatedly sampling its values from the dataset with replacement. It's highly
versatile since it doesn't rely on strong distributional assumptions.


\begin{center}
\begin{figure}[ht]
  \centering
  \includegraphics[width=5in]{img/bootstrap.jpg}
  \caption{Illustration of bootstrap resampling: The distribution of the
  statistic (e.g., mean) of the target variable for two different samples
(called groups). This demonstrates the variability of the statistic across
resampled datasets.}
  \label{Figure:fig_beh}
\end{figure}
\end{center}

\chapter{Exploratory Data Analysis}
\label{cha:eda}
%---------------------------------------------------------------------------


\section{Spotify Features}
\label{sec:spotifyfeatures}

\subsection*{Correlation Heatmap}
\label{sec:correlationheatmapsspotifyfeatures}

\begin{center}
\begin{figure}[H]
  \centering
  \includegraphics[width=4in]{img/corr_heatmap_spotify_features.png}
  \caption{Pearson's Correlation Heatmap of Spotify Features}
  \label{Figure:fig_beh}
\end{figure}
\end{center}

\subsection*{Hierarchical Clustering}
\label{sec:hierarchicalclustering}

\begin{center}
\begin{figure}[H]
  \centering
  \includegraphics[width=4in]{img/dendrogram_spotify_features.png}
  \caption{Hierarchical Clustering of Spotify Features}
  \label{Figure:dendrogram_spotify_features}
\end{figure}
\end{center}

\subsection*{Observations}
Spotify audio features show high level of correlation between each other, especially:
\begin{itemize}
  \item \textit{Energy} and \textit{Loudness}: A high correlation(0.74) suggest
    that energetic songs are typically louder.
  \item \textit{Danceability} and \textit{Valence}: high correlation between
    them indicates that songs perceived as positive and happy are usually more
    danceable.
  \item \textit{Energy} and \textit{Acousticness}: there is a strong negative
    correlation between those features, suggesting that high-energy tracks are
    less likely to have acoustic elements.
  \item The dendrogram\ref{Figure:dendrogram_spotify_features} shows relations
    between features and allows us to compare which features correlate with
    each other. In addition to the relationships seen in the correlation
    heatmap, we observe that \textit{popularity} appears to have a weak
    correlation with \textit{release year} and \textit{speechiness}.

\end{itemize}


%---------------------------------------------------------------------------


\section{Lyrical Features}

\subsection*{Correlation Heatmap}
\label{sec:correlationheatmapsspotifyfeatures}

\begin{center}
\begin{figure}[H]
  \centering
  \includegraphics[width=4in]{img/corr_heatmap_lyrical.png}
  \caption{Pearson's Correlation Heatmap of Lyrical Features}
  \label{Figure:fig_beh}
\end{figure}
\end{center}

\subsection*{Hierarchical Clustering}
\label{sec:hierarchicalclustering}

\begin{center}
\begin{figure}[H]
  \centering
  \includegraphics[width=4in]{img/dendrogram_lyrical.png}
  \caption{Hierarchical Clustering of Lyrical Features}
  \label{Figure:dendrogram_spotify_features}
\end{figure}
\end{center}


\subsection*{Observations}

\begin{itemize}
  \item \textit{Flesch Reading Ease}, \textit{Gunning Fog} and \textit{Dale
    Chall} scores exhibit strong positive correlations, highlighting their
    shared focus on measuring lyrical complexity.
  \item \textit{Lexical Richness} correlates with \textit{Word Count},
    indicating that more lexically rrich lyrics tend to have  greater variety
    of words.
  \item On the dendrogram we can distinguish three major clusters of features:
    \begin{itemize}
      \item \textbf{Lyrical Complexity Metrics} - \textit{Flesch Reading Ease},
        \textit{Gunning Fog}, \textit{Dale Chall} and \textit{Syntactic
        Complexity } quantify how difficult and complex the lyrics are.
      \item  \textbf{Lexical Features} - features such as \textit{Type-Token
        Ratio}, \textit{Lexical Richness} and \textit{Unique Word Count} form a
        cohesive cluster.
      \item \textbf{Sentiment-Related Features} - it includes \textit{Sentiment
        Polarity}, \textit{VADER Compound} and \textit{Semantic Depth}, which
        reflect emotional aspects of the lyrics.
    \end{itemize}
  \item Sentiment-related features are relatively closely grouped, suggesting a
    high level of interdependence
\end{itemize}


%---------------------------------------------------------------------------
\section{Audio Features}

\subsection*{Correlation Heatmap}
\label{sec:correlationheatmapsspotifyfeatures}

\begin{center}
\begin{figure}[H]
  \centering
  \includegraphics[width=4in]{img/corr_heatmap_audio.png}
  \caption{Pearson's Correlation Heatmap of Audio Features}
  \label{Figure:fig_beh}
\end{figure}
\end{center}

\subsection*{Hierarchical Clustering}
\label{sec:hierarchicalclustering}

\begin{center}
\begin{figure}[H]
  \centering
  \includegraphics[width=4in]{img/dendrogram_audio.png}
  \caption{Hierarchical Clustering of Audio Features}
  \label{Figure:dendrogram_spotify_features}
\end{figure}
\end{center}


\subsection*{Observations}
\begin{itemize}
  \item \textit{MFCC Features} show strong correlations among themselves(e.g.
    \textit{mfcc\_4}, \textit{mfcc\_6}, \textit{mfcc\_8}, etc.).
  \item Similarily \textit{chroma} features are highly correlated with each
    other, indicating that  they capture similar aspects of tonal energy
    distribution.
  \item \textit{spectral\_contrast} features exhibit relatively small
    correlations with \textit{mfcc} and  \textit{chroma} features, suggesting
    that they capture different characteristics of audio
  \item All \textit{chroma} features are grouped together in the same cluster
    on the dendrogram, supporting the conclusion drawn from the heatmap about
    strong correlations between them
  \item \textit{zcr} and \textit{tempo\_extracted} are relatively independent
  \item A certain degree of overlap between \textit{mfcc} and
    \textit{spectral\_contrast} features is observed in the dendrogram,
    indicating potential shared information in capturing specific audio
    properties.
\end{itemize}

%---------------------------------------------------------------------------


\section{Empath Features}

\subsection*{Correlation Heatmap}
\label{sec:correlationheatmapsspotifyfeatures}

\begin{center}
\begin{figure}[H]
  \centering
  \includegraphics[width=4in]{img/corr_heatmap_empath.png}
  \caption{Pearson's Correlation Heatmap of Empath Features}
  \label{Figure:fig_beh}
\end{figure}
\end{center}

\subsection*{Hierarchical Clustering}
\label{sec:hierarchicalclustering}

\begin{center}
\begin{figure}[H]
  \centering
  \includegraphics[width=4in]{img/dendrogram_empath.png}
  \caption{Hierarchical Clustering of Empath Features}
  \label{Figure:dendrogram_spotify_features}
\end{figure}
\end{center}


\subsection*{Observations}

\begin{itemize}
  \item \textbf{High Correlations Reflect Logical Groupings}: high correlation
    between features such as \textit{empath\_optimism}, \textit{empath\_love} and
    \textit{empath\_affection} align with the intuitive understanding that
    these aspects are closely related.
  \item Observed patters in corrrelations confirm that Empath was designed to
    group semantically related concepts together.
  \item The visualizations align with the intended design of Empath as a tool
    for interpretable feature  extraction.
\end{itemize}


%---------------------------------------------------------------------------

\section{Genre Analysis}

\subsection{Genre Similarity}

\begin{center}
\begin{figure}[H]
  \centering
  \includegraphics[width=5in]{img/genres_dendrogram.png}
  \caption{Hierarchical Clustering of Genres by lyrical and audio features}
  \label{Figure:dendrogram_spotify_features}
\end{figure}
\end{center}

\begin{center}
\begin{figure}[H]
  \centering
  \includegraphics[width=4in]{img/genres_similarity_heatmap.png}
  \caption{Heatmap of euclidean distance of Genres calculated on lyrical and audio features}
  \label{Figure:dendrogram_spotify_features}
\end{figure}
\end{center}

The dendrogram illustrates hierarchical clustering of genres based on their
lyrical and  audio features. The clustering aligns with general cultural and
musical understandings of  these genres. For instance, closely related  genres
like hip hop and rap are grouped  together, indicating that they're highly
similar.

The heatmap further supports this observation. The Euclidean distances of
\textit{rap} and \textit{hip hop} with other genres stand  out as the highest.

Interestingly, genres such as \textit{metal} and \textit{indie}, while distinct
in sound, share some overlap in features, which is reflected on the dendrogram.

\begin{center}
\begin{figure}[H]
  \centering
  \includegraphics[width=6in]{img/heatmap_of_empath.png}
  \caption{Heatmap of mean values of empath features in each genre}
  \label{Figure:dendrogram_spotify_features}
\end{figure}
\end{center}

The heatmap displays scaled mean values of \textit{Empath features} for
different musical genres, highlighting thematic  differences in their lyrics.
It can be observed that:
\begin{itemize}
  \item \textbf{Country} music lyrics often involve topics related to family,
    love and  positive emotions. It seems to align well with genre's tendency
    to narrate personal, heartfelt and often nostalgic stories.
  \item \textbf{EDM} lyrics seem to often touch upon love and optimism.
  \item For both \textbf{Hip Hop} and \textbf{Rap} according to Empath  the
    most prominent topic is \textit{speaking}.
  \item For \textbf{Metal} High values in \textit{negative emotion }and
    \textit{pain} underscore the genre's focus on intense, darker themes.
  \item Elevated values in \textit{love} and \textit{positive emotion}
    highlight \textbf{Pop} music’s focus on themes of love and relationships.
  \item High values in \textit{friends } and \textit{optimism} categories
    reflect \textbf{Reggae's} emphasis on positivity, social connection, and
    uplifting messages.
\end{itemize}

\subsection{Lyrical Similarity Based on Embeddings}
\begin{center}
\begin{figure}[H]
  \centering
  \includegraphics[width=6in]{img/tsne_genres.png}
  \caption{Heatmap highlighting the maximum feature values across genres.}
  \label{Figure:dendrogram_spotify_features}
\end{figure}
\end{center}

\textit{t-SNE (t-Distributed Stochastic Neighbor Embedding)} was applied to the
Euclidean distance matrix of mean standardized Word2Vec and TF-IDF embeddings
for each genre. This dimensionality reduction technique allows to project
high-dimensional data into a 2D space while preserving the relative distances
and similarities between data points as much as possible. This enables to
visualize lyrical similarities between different genres.

In order to further explore lyrical relationships between the genres K-means
clustering was applied on the mean standardized embeddings per genre,The
optimal number of clusters was determined using the elbow method, which
suggested four distinct clusters. This additional operation clusters the most
similar genres together.

\begin{itemize}
  \item \textit{Hip hop} and \textit{rap} form a distinct cluster, indicating
    strong lyrical and thematic similarity.
  \item \textit{Reggae} is slightly separate from other genres but still part
    of a larger cluster that includes \textit{rock}, \textit{pop},
    \textit{country}, and \textit{R\&B}, suggesting moderate similarity in
    lyrical and thematic content of songs belonging to those genres.
  \item \textit{Metal} and \textit{indie} are closely positioned, sharing
    overlapping themes while forming a separate cluster.
  \item \textit{EDM} is positioned furthest from other genres, highlighting
    its unique lyrical and thematic style.
\end{itemize}

\subsection{Top Genre Characteristics}
\begin{center}
\begin{figure}[H]
  \centering
  \includegraphics[width=6in]{img/heatmap_max_feature_values_by_genre.png}
  \caption{Heatmap highlighting the maximum feature values across genres.}
  \label{Figure:dendrogram_spotify_features}
\end{figure}
\end{center}
Each filled cell represents the genre that exhibits the highest mean value for
the corresponding feature, calculated using scaled feature values. This
visualization emphasizes distinctive characteristics of each genre, and allows
to find out ``in  which genre the feature achieved its highest mean``.

\begin{itemize}
  \item \textbf{Country}: Highest in acousticness and rhyme density
  \item \textbf{EDM}: Dominated in tempo, loudness, and featured the most
    recent songs
  \item \textbf{Hip Hop}: Excelled in unique word count, lexical richness, and
    semantic depth.
  \item \textbf{Metal}: Stood out for energy and longer song durations.
  \item \textbf{Pop}: Showed the highest positivity (VADER compound) and
    subjectivity in lyrics.
  \item \textbf{Rap}: Highest in repetition count and reading ease.
  \item \textbf{Reggae}: Highlighted by high valence and danceability.
  \item \textbf{Rock}: Scored highest in in popularity.
    
\end{itemize}

\begin{center}
\begin{figure}[H]
  \centering
  \includegraphics[width=6in]{img/heatmap_top_feature_values_by_genre.png}
  \caption{Heatmap highlighting the top 5 features by genre. Each filled cell
  represents one of the five highest mean feature values for a given genre,
calculated on scaled data. This visualization focuses on identifying the most
distinctive traits for each genre.}
  \label{Figure:dendrogram_spotify_features}
\end{figure}
\end{center}

Lastly, we switch perspectives, focusing not on individual features but instead
examining each genrer to identify the top 5 features that characterize it. Each
filled cell in that heatmap shows one of the five most prominent
characteristics exhibited by that genre.


\chapter{Experiments and Results}
\label{cha:experimentsandresults}
%---------------------------------------------------------------------------

\section{Song Popularity}
\label{sec:songpopularity}

The analysis of song popularity provides valuable insights into the factors
that influence the success of music tracks. This problem has two approaches:
\begin{itemize}
  \item \textbf{Regression} - Spotify's popularity is a value on scale of
    0-100. This approach involves training a regression model trying to predict
    that value.
  \item \textbf{Classification} - assigning binary label to the songs(popular
    vs. unpopular) and training a classification model to predict it.
\end{itemize}

In this section the prediction of popularity was attempted using several
regression and classification models. These models were trained on different
sets of features, including Spotify metadata, lyrical attributes, and audio
features, with the aim of investigating the predictive power of those features
and their impact on popularity.

Catboost models were used as the primary predictive tool due to their
robustness and performance in handling complex relationships. 

Baseline models were also implemented to serve as a reference point:
\begin{itemize}
  \item \textit{Baseline Mean Model}: baseline model for regression that always predicts the mean.
  \item \textit{Baseline Majority Model}: baseline model for classification that always predicts the majority class.
  \item \textit{Baseline Random Model}: baseline model for classification that predicts random class.
\end{itemize}

The experiments involved both quantitative evaluatoin and SHAP analysis to
assess feature importance and interpretability of the models.

%---------------------------------------------------------------------------
\subsection{Regression Approach}

\begin{center}
\begin{figure}[H]
  \centering
  \includegraphics[width=6in]{img/reg_pipeline.png}
  \caption{Regression model pipeline. It involves preprocessing, feature
  selection and CatBoost model.}
  \label{Figure:fig_beh}
\end{figure}
\end{center}

The pipeline was fit with all available features, that includes Spotify audio
features and metadata, lyrical features and features extracted from the audio
files. The feature selection step chose the subset of most valuable features
based on the feature importance from initial Random Forest model.

\begin{table}[H]
\centering
\caption{Results of regression of popularity.}
\begin{tblr}{
  hline{2} = {-}{},
}
 & \textbf{Model}      & \textbf{Features} & \textbf{MAE}  & \textbf{RMSE} & \textbf{$R^2$} \\
 & Baseline Mean Model &                   & 13.53         & 16.90         & -0.01          \\
 & Catboost            & all               & \textbf{6.33} & \textbf{8.52} & \textbf{0.74}  \\
 & Catboost            & lyrical           & 11.81         & 15.29         & 0.17           \\
 & Catboost            & spotify data      & \textbf{6.04} & \textbf{8.10} & \textbf{0.77}  \\
 & Catboost            & audio             & 12.71         & 16.15         & 0.08           
\end{tblr}
\end{table}

As seen in the results table, the CatBoost model trained on spotify data only
achieved the best performance, with \textbf{Mean Absolute Error(MAE) of
\textbf{6.04} and and $R^2$ of 0.77}, closely followed by the model trained on
all features that achieved \textbf{MAE of 6.33 and ($R^2$ of 0.74}. The results
clearly indicate the importance of Spotify audio features and metadata in
prediction of song's popularity. Predicting the success  of a musical track
based on solely the lyrics or acoustic features turned out to be diffifcult and
models trained on those subsets of features showed only slight improvement in
comparison to the baseline model. In contrast, the model trained on Spotify data only
achieved \textbf{55.3\%} better MAE score than the baseline model.

Despite feature selection, the model trained on all features performed slightly
worse than the model trained on Spotify features alone. This likely occurred
because adding less important or redundant features reduced the model's ability
to focus on the most relevant Spotify features. The additional features may
have introduced noise, making it harder to identify clear patterns. This
behavior might also hint at slight overfitting on the redundant features,
despite the use of cross-validation. Increased regularization could potentially
mitigate this issue in future experiments.


SHAP analysis was conducted to interpret the predictions of the regression
model and asses the importance of individual features and their contribution to
song's popularity.


\begin{center}
\begin{figure}[H]
  \centering
  \includegraphics[width=5in]{img/feature_importance_popularity_reg.png}
  \caption{SHAP feature importance plot of the regression model for popularity trained on all features.}
  \label{Figure:fig_beh}
\end{figure}
\end{center}


The mean absolute SHAP values were used to rank the overall importance of input
features. Key insights:
\begin{itemize}
  \item \textbf{Dominance of \textit{Genre} and \textit{Album Release Year}.}:
    these two features account for the majority of the predictive power in the
    model. Genre captures the musical style and release year reflects trends
    and cultural preferences over time.
  \item Spotify's audio features like \textit{loudness} and
    \textit{speechiness} showed moderate contribution to model's performance.
    Their score highlights their relevance in describing popular tracks' audio
    characteristics.
  \item Lyric-based features like \textit{word count} and \textit{unique word
    count} were significantly less impactful, however still contributed to
    model's performance to some degree.
\end{itemize}

\begin{center}
\begin{figure}[H]
  \centering
  \includegraphics[width=5in]{img/beeswarm_popularity_reg.png}
  \caption{SHAP beeswarm plot of the regression model for popularity trained on all features.}
  \label{Figure:fig_beh}
\end{figure}
\end{center}

\begin{itemize}
  \item Higher values of \textit{Release Year} strongly correlate with higher
    \textit{Popularity}, indicating that newer songs are generally more
    popular.
  \item Higher \textit{loudness} and lower \textit{speechiness}
    were associated with increased \textit{popularity}, reflecting their
    influence on listener preferences and music trends.
\end{itemize}



%---------------------------------------------------------------------------
\subsection{Classification Approach}

The task of predicting song popularity was reformulated as a binary
classification problem, where the target was to determine whether a song is
"popular" (1) or "not popular" (0). In order to create those labels from an
integer variable with range 0-100, the \textbf{70th perrcentile of the
popularity values was used as the threshold}. Songs with popularity greater or
equal  than this  threshold werre  labeled as popular, and the rest was labeled
as unpopular. This thresholding approach based on quantile ensures a balanced
representation of popular songs in the dataset while accounting for the
naturally skewed distribution of  popularity scores.

% \usepackage{tabularray}
\begin{table}[H]
\centering
\caption{Results of classification of popularity.}
\begin{tblr}{
  hline{2} = {-}{},
}
 & \textbf{Model}          & \textbf{Features} & \textbf{Accuracy} & \textbf{F1(w.avg.)} \\
 & Baseline Majority Model &                   & 65.93\%           & 52.39\%             \\
 & Baseline Random Model   &                   & 48.27\%           & 49.57\%             \\
 & Catboost                & all               & \textbf{84.68\%}  & \textbf{84.71\%}    \\
 & Catboost                & lyrical           & 65.79\%           & 64.77\%             \\
 & Catboost                & spotify data      & \textbf{82.06\%}  & \textbf{82.32\%}    \\
 & Catboost                & audio             & 63.17\%           & 63.00\%             
\end{tblr}
\end{table}


The classification models were evaluated using accuracy and weighted F1-score.
\textit{Baseline Majority Model} (which predicts the majority class for all
samples) achieved an accuracy of 65.93\% and an F1-score of 52.39\%. This
reflects the class imbalance introduced by the thresholding, with a higher
proportion of songs being labeled as "not popular."

The \textit{Baseline Random Model}, which guesses classes randomly, performed
worse in accuracy (48.27\%) and achieved a slightly higher F1-score of 49.57\%
compared to accuracy due to its balanced attempts to handle both classes.

The CatBoost model significantly outperformed both baselines, achieving an
overall accuracy of 84.68\% and a weighted F1-score of 84.71\% when trained on
all features.This result demonstrates the effectiveness of the full feature set
in capturing the complexity of popularity prediction.


The CatBoost model trained on Spotify data performed well, with an accuracy of
82.06\% and an F1-score of 82.32\%, showing that Spotify features alone provide
strong predictive power. However, the model trained on all features slightly
surpassed it, suggesting that additional features offer some complementary
information.

The model using only lyrical features achieved an accuracy of 65.79\%, lower
than the \textit{Baseline Majority Model}, but a higher F1-score of 64.77\%,
reflecting its ability to handle class imbalance better than the majority
model. This confirms that lyrical features are limited in their ability to
predict song popularity.


The performance of the model trained on features extracted from audio files was
similar to the lyrical one, demonstrating that isolated audio features are not
sufficient to fully predict popularity.


SHAP analysis revealed that key features such as \textit{genre}, \textit{album
release year} and \textit{loudness} were the most
significant predictors of popularity. These findings were consistent with the
regression task, reinforcing the robustness of these features across different
modeling approaches. This consistency highlights their central role in
understanding and predicting song success.


\begin{center}
\begin{figure}[H]
  \centering
  \includegraphics[width=5in]{img/feature_importance_popularity_clf.png}
  \caption{SHAP feature importance plot of the classification model for popularity trained on all features.}
  \label{Figure:fig_beh}
\end{figure}
\end{center}

\begin{center}
\begin{figure}[H]
  \centering
  \includegraphics[width=5in]{img/beeswarm_popularity_clf.png}
  \caption{SHAP beeswarm plot of the classification model for popularity trained on all features.}
  \label{Figure:fig_beh}
\end{figure}
\end{center}



The SHAP analysis for the classification model highlights similar dependencies
to the regression model. Key features like \textit{genre} and
\textit{album release year} remain the most important, emphasizing their role
in shaping song popularity. 

Several differences can be observed between feature importance in the
regression and classification model, notably:
\begin{itemize}
  \item Lyrics embeddings(TF-IDF and Word2Vec) seemed to play a bigger role in
    the classification model.
  \item Classification model seemed to pay much less attention to Spotify audio
    features, like \textit{danceability}, \textit{energy} and
    \textit{speechiness}.
  \item Unlike in regression, in classification a sentiment metric,
    \textit{VADER compound} contributed to the prediction of popularity. More
    popular songs tend to have more positive lyrics.
\end{itemize}


%---------------------------------------------------------------------------
\section{Explicitness}
\label{sec:explicitness}

\subsection{Classification Approach}
The task of predicting whether a song contains explicit content was approached
as a classification problem. The performance of the CatBoost model was compared
against baseline on different feature subsets. One of the key challenges of
this problem was significant class imbalance present in the dataset;
approximately 85\% of the songs did not contain explicit content.

Significant class imbalance can lead to models favoring the majority class,
potentially resulting in  high accuracy but poor performance on the minority
class. To address this problem, CatBoost's class weights parameter was used.
This parameter allowed the model to penalize misclassifications of the minority
class more heavily, therefore improving its ability to recognize explicit
content.

% \usepackage{tabularray}
\begin{table}[H]
\centering
\caption{Results of classification of explicitness.}
\begin{tblr}{
  hline{2} = {-}{},
}
 & \textbf{Model}          & \textbf{Features} & \textbf{Accuracy} & \textbf{F1 weighted average} \\
 & Baseline Majority Model &                   & 84.00\%           & 76.69\%                      \\
 & Baseline Random Model   &                   & 47.03\%           & 53.91\%                      \\
 & Catboost                & all               & \textbf{92.41\%}  & \textbf{92.40\%}             \\
 & Catboost                & lyrical           & \textbf{92.41\%}  & \textbf{92.25\%}             \\
 & Catboost                & spotify data      & 86.48\%           & 87.11\%                      \\
 & Catboost                & audio             & 82.75\%           & 82.47\%                      
\end{tblr}
\end{table}


The table presents the results of the models in terms of accurracy and weighted
F1 score. The \textit{Baseline Majority Model} achieved accuracy of 84\% and
significantly lower F1 weighted average of 76.69\%, which reflects its inability
to handle class balance effectively. The \textit{Baseline Random Model}
performed very poorly, with an accuracy of 47.03\% and F1 weighted average of
53.91\%.

In contrast, the CatBoost model outperformed both baselines achieving accuracy
of \textbf{92.41\%} and weighted average score of \textbf{92.4\%} when trained
on all features. Interestingly, the model trained only using lyrical features
performed nearly as well, indicating that explicitness can largely be predicted
based on lyrical content. Models that relied on Spotify metadata and features
extracted from audio files had lower performance, achieving accurracy close to
\textit{Baseline Majority Model}, but with higher F1 scores, due to their
ability to address class imbalance. These observations emphasize the centrality
of lyrical information in prediction of explicit content.


\begin{center}
\begin{figure}[H]
  \centering
  \includegraphics[width=5in]{img/beeswarm_explicitness.png}
  \caption{SHAP beeswarm plot of the classification model for explicitness trained on all features.}
  \label{Figure:fig_beh}
\end{figure}
\end{center}

\begin{center}
\begin{figure}[H]
  \centering
  \includegraphics[width=5in]{img/feature_importance_explicitness.png}
  \caption{SHAP beeswarm plot of the classification model for explicitness  trained on all features.}
  \label{Figure:fig_beh}
\end{figure}
\end{center}


Based on the SHAP analysis we can observe that the key features for this task
turned out to be:
\begin{itemize}
  \item \textbf{Album Release Year}: this feature had  highest feature
    importance, reflecting temporal trends in explicit content. Newer songs are
    statistically more likely to contain explicit language, reflecting changing
    societal norms and artistic expressions over time.
  \item \textbf{Empath Swearing Terms}: As expected, the presence of swearing
    terms was a strong indicator of explicitness.
  \item \textbf{Genre} was the third most important feature, showing that
    certain genres  are more likely to contain explicit content than others.
  \item \textbf{TF-IDF}: several vectors created using TF-IDF and later reduced
    via PCA significantly contributed to the model's predictions. Similar to
    the Empath feature, TF-IDF likely captured swear words and related patterns
    in the lyrics, reinforcing its importance.
  \item \textbf{Speechiness}: songs with higher values of
    speechiness—indicating a greater presence of spoken-word elements—were more
    likely to be labeled as explicit.
\end{itemize}




The task of predicting explicitness could potentially achieve higher
performance by leveraging the interpretability of explicitness, which provides
clear hints about the most relevant features for this property. However, the
primary goal of this analysis was not to maximize model performance but to
explore the impact of various feature subsets on the prediction of
explicitness.

The results indicate that predicting explicitness based solely on audio
features is inherently challenging. The model trained on audio features
performed worse than the baseline, suggesting that the extracted audio
characteristics from the MP3 files do not sufficiently capture the explicit
nature of songs. This highlights the need for features more directly related to
lyrical or contextual information when modeling explicit content.



%---------------------------------------------------------------------------
\subsection{Impact of Explicit Language on Popularity and Sentiment}
\label{sec:explicitmorepopular}

The relationship between explicit content and song popularity is a topic of
interest in understanding the cultural and commercial dynamics of music.
Explicit songs often reflect bold themes, which may resonate more strongly with
certain audiences. 

To investigate this, a bootstrap test was conducted to determine whether
explicit songs are, on average, more popular than non-explicit songs. Bootstrap
was chosen for this analysis due to its robustness and flexibility. Unlike
traditional parametric tests, the bootstrap method does not rely on strong
assumptions about the underlying data distribution, making it well-suited for
analyzing real-world datasets that may not meet normality or other strict
requirements.


\begin{center}
\begin{figure}[H]
  \centering
  \includegraphics[width=5in]{img/explicitness_bootstrap.png}
  \caption{Bootstrap Test to check if explicit songs are on average more
  popular.}
  \label{Figure:fig_bh}
\end{figure}
\end{center}



\begin{center}
\begin{figure}[H]
  \centering
  \includegraphics[width=5in]{img/explicitness_bootstrap_2.png}
  \caption{Bootstrap Test to check if explicit songs are on average more
  positive or negative than non-explicit songs.}
  \label{Figure:fig_bh}
\end{figure}
\end{center}



\begin{table}[H]
\centering
\caption{Results of the Bootstrap Test.}
\resizebox{\linewidth}{!}{%
\begin{tblr}{
  cell{2}{4} = {fg=green},
  cell{3}{4} = {fg=red},
  hlines,
  vlines,
}
\textbf{Metric}     & \textbf{Control Group} & \textbf{Treatment Group} & \textbf{pct\_diff} & \textbf{pct\_ci\_lower} & \textbf{\textbf{pct\_ci\_upper}} & \textbf{abs\_diff} & \textbf{abs\_diff\_ci\_lower} & \textbf{abs\_diff\_ci\_upper} \\
Popularity          & non-explicit           & explicit                 & \textbf{2.96\%}    & 0.08\%                  & 5.79\%                           & 1.74               & 0.0453                        & 3.4013                        \\
Sentiment~ Polarity & non-explicit           & explicit                 & \textbf{-87.10\%}  & -100.69\%               & -73.63\%                         & -0.09              & -0.1067                       & -0.0780                       
\end{tblr}
}
\end{table}


The analysis reveals that explicit songs are, on average, \textbf{2.96\% more
popular} than non-explicit songs. The confidence interval of [0.08\%, 5.79\%]
suggests statistically significant but modest increase in popularity for
explicit songs.

For sentiment polarity, explicit songs turned out to be on average
\textbf{87.10\% less positive} than non-explicit songs. The confidence interval
of [-100.69\%, -73.63\%] suggests a very significant difference. This
observation aligns with the notion that bold songs that contain explicit
language often reflect intense or provocative themes, which may carry less
positive emotion.

%---------------------------------------------------------------------------
\section{Sentiment}
\label{sec:sentiment}
\begin{center}
\begin{figure}[H]
  \centering
  \includegraphics[width=7in]{img/wordclouds.png}
  \caption{Wordclouds of Positive and Negative Sentiment Songs.}
  \label{Figure:fig_bh}
\end{figure}
\end{center}

This section tackles the problem of song lyrics sentiment prediction. The
binary sentiment labels(positive (1) and negative (0)) were derived from the
variable \textit{sentiment polarity}. Songs with polarity values below 0 were
labeled as negative, while those above 0 were labeled as positive.

This process resulted in an unbalanced target variable, where the positive
class was overrepresented. Additionally, due to the left-skewed distribution of
\textit{sentiment polarity}, songs labeled as \textit{positive} were, on
average, more strongly positive compared to the relatively moderate negativity
of songs in the \textit{negative} class.

While this imbalance and class definition might pose challenges for traditional
predictive modeling, it is not a major concern for this study. The primary
objective is to understand which factors contribute to a song's sentiment
rather than achieving perfect predictive accuracy. Chosen approach
provides sufficient insight into the relationship between song features and
sentiment.



\begin{table}[H]
\centering
\caption{Results of classification of sentiment.}
\begin{tblr}{
  hline{2} = {-}{},
}
 & \textbf{Model}          & \textbf{Features} & \textbf{Accuracy} & \textbf{F1(w.avg.)} \\
 & Baseline Majority Model &                   & 69.65\%           & 57.19\%             \\
 & Baseline Random Model   &                   & 52.55\%           & 54.46\%             \\
 & Catboost                & all               & 71.44\%           & 69.77\%             \\
 & Catboost                & lyrical           & \textbf{71.86\%}  & \textbf{70.72\%}    \\
 & Catboost                & spotify data      & 67.44\%           & 64.94\%             \\
 & Catboost                & audio             & 65.10\%           & 63.43\%             
\end{tblr}
\end{table}




The \textit{Baseline Majority Model}, which simply predicted the most common
class for all samples, achieved an accuracy of 69.65\% and an F1-score of
57.19\%, reflecting the class imbalance in the dataset. The \textit{Baseline
Random Model}, which predicted classes randomly, performed worse in terms of
accuracy (52.55\%) but slightly better in F1-score (54.46\%).

Among the CatBoost models, the one trained solely on lyrical features performed
the best, achieving an accuracy of 71.86\% and a weighted F1-score of 70.72\%.
This highlights the importance of lyrical features in capturing sentiment. The
model trained on all features followed closely, but performed slightly worse
due to large amount of redundant features, despite the presence of feature
selection step in the pipeline that attempted to discard them.

Models trained on only Spotify or audio features performed worse, achieving
accuracies of 67.44\% and 65.10\%, respectively. These results suggest that
while Spotify and audio features capture useful information, they are less
indicative of sentiment compared to lyrical data. 

The results indicate that features extracted directly from MP3 files are not
strong predictors of a song's sentiment, despite the intuitive assumption that
audio characteristics should play a role. This outcome suggests that the
current set of extracted audio features may not sufficiently capture the
nuanced aspects of sentiment. Future work could focus on advanced feature
engineering or exploring additional audio-based descriptors to improve the
performance of audio-based sentiment prediction models.

\begin{center}
\begin{figure}[H]
  \centering
  \includegraphics[width=5in]{img/beeswarm_sentiment.png}
  \caption{SHAP beeswarm plot of the classification model for sentiment
  trained on all features.}
  \label{Figure:fig_bh}
\end{figure}
\end{center}

\begin{center}
\begin{figure}[H]
  \centering
  \includegraphics[width=5in]{img/feature_importance_sentiment.png}
  \caption{SHAP feature importance plot of the classification model for
  sentiment trained on all features.}
  \label{Figure:fig_eh}
\end{figure}
\end{center}

The SHAP analysis showed that lyrical features, like \textit{empath negative
emotion}, \textit{empath optimism}, and \textit{empath swearing terms}, were
the most influential in predicting sentiment, with negative emotions and
swearing linked to negative sentiment, and optimism and positive emotion tied
to positive sentiment. TF-IDF and Word2Vec embeddings also contributed,
reflecting patterns in lyrics. Audio features, however, had minimal impact,
emphasizing the need for improved audio feature engineering to capture
sentiment.

%---------------------------------------------------------------------------
\section{Genre}
\label{sec:genre}

This task aimed to predict song's genre using  various  feature subsets. The
dataset included 10 distinct genres, covering a diverse range of musical
styles.

\begin{table}[H]
\centering
\caption{Results of classification of genre.}
\resizebox{\linewidth}{!}{%
\begin{tblr}{
  hline{2} = {-}{},
}
 & \textbf{Model}          & \textbf{Features} & \textbf{Accuracy} & \textbf{F1(w.avg.)} \\
 & Baseline Majority Model &                   & 11.17\%           & 02.24\%             \\
 & Baseline Random Model   &                   & 10.20\%           & 10.38\%             \\
 & Catboost                & all               & \textbf{61.24\%}  & \textbf{60.27\%}    \\
 & Catboost                & lyrical           & 32.96\%           & 32.23\%             \\
 & Catboost                & spotify data      & 56.96\%           & 56.65\%             \\
 & Catboost                & audio             & 31.31\%           & 30.33\%             
\end{tblr}
}
\end{table}

The baseline models achieved very poor performance. The \textit{Baseline
Majority Model} achieved accuracy of 11.17\% and weighted F1 of 2.24\%, and the
\textit{Baseline Random Model} accuracy of 10.20\% and F1 of 10.38\%.The
CatBoost model trained on all features significantly outperformed the
baselines, achieving 61.24\% accuracy and 60.27\% weighted F1-score. This
highlights the importance of combining diverse feature sets (lyrics, Spotify
data, audio features, embeddings) for genre prediction.

Among the individual feature subsets, Spotify features turned out to be the
most predictive, resulting in a relatively well-performing model. While lyrical
and audio features alone were not strong predictors of genre, their combination
with Spotify features significantly improved the model's performance,
highlighting the complementary nature of these data sources.





\begin{center}
\begin{figure}[H]
  \centering
  \includegraphics[width=5in]{img/feature_importance_genre_clf.png}
  \caption{SHAP feature importance plot of the classification model for genre
  trained on all features. The scores were averaged across all classes.}
  \label{Figure:fig_eh}
\end{figure}
\end{center}

Looking at the SHAP global feature importance plot(values averaged across all
classes) we can observe that:
\begin{itemize}
  \item \textit{Popularity} and \textit{album release year} emerged as the most
    critical predictors of \textit{genre}. These two features likely capture
    trends and listener preferences linked to specific genres.
  \item Other key features included \textit{danceability},
    \textit{speechiness}, and \textit{unique word count}, which relate to the
    style of the music and the content of the lyrics, helping to distinguish
    between genres. \item Audio features like \textit{acousticness} and
    \textit{spectral contrast} also played a significant role, highlighting
    that the sound of the music reflects its genre. 
  \item The \textit{explicit} flag also had a noticeable effect, which makes
    sense given its strong connection to genres like rap and hip-hop.

\end{itemize}

%---------------------------------------------------------------------------


\section{Topics Modelling}
\label{sec:topicsmodelling}

\subsection{Latent Dirichlet Allocation}

\begin{center}
\begin{figure}[H]
  \centering
  \includegraphics[width=4in]{img/topics/png/topics.png}
  \caption{Intertopic Distance Map.}
  \label{Figure:intertopicdistancemap}
\end{figure}
\end{center}

\begin{itemize}
  \item The bubble chart shows the intertopic distance map using multidimensional
scaling (MDS). Each bubble represents a topic extracted by the LDA model.
  \item The position of the bubbles reflects the relationship between topics.
    Topics closer together share more words or themes in common.
  \item The size of the bubbles represents the prevalence of the topic within
    the entire dataset. Larger bubbles correspond to topics with a higher
    proportion of terms in the dataset.
\end{itemize}



\begin{center}
\begin{figure}[H]
  \centering
  \includegraphics[width=3in]{img/topics/png/general.png}
  \caption{Top-30 Most Salient Terms across the entire corpus. The bars
  represent how often each term appears in the entire corpus.}
  \label{Figure:fig_eh}
\end{figure}
\end{center}

 \textbf{Blue bars show overall frequency of the terms in the entire corpus.
 Red bars indicate estimated frequency of the terms in that topic.} Now we can
 look at the results in each topic:

\begin{center}
\begin{figure}[H]
  \centering
  \includegraphics[width=3in]{img/topics/png/t1.png}
  \caption{Top-30 Most Relevant Terms in Topic 1.}
  \label{Figure:t1}
\end{figure}
\end{center}

\begin{center}
\begin{figure}[H]
  \centering
  \includegraphics[width=3in]{img/topics/png/t2.png}
  \caption{Top-30 Most Relevant Terms in Topic 2.}
  \label{Figure:fig_eh}
\end{figure}
\end{center}

\begin{center}
\begin{figure}[H]
  \centering
  \includegraphics[width=3in]{img/topics/png/t3.png}
  \caption{Top-30 Most Relevant Terms in Topic 3.}
  \label{Figure:t3}
\end{figure}
\end{center}

\begin{center}
\begin{figure}[H]
  \centering
  \includegraphics[width=3in]{img/topics/png/t4.png}
  \caption{Top-30 Most Relevant Terms in Topic 4.}
  \label{Figure:fig_eh}
\end{figure}
\end{center}

\begin{center}
\begin{figure}[H]
  \centering
  \includegraphics[width=3in]{img/topics/png/t5.png}
  \caption{Top-30 Most Relevant Terms in Topic 5.}
  \label{Figure:fig_eh}
\end{figure}
\end{center}


The topic modeling results reveal distinct themes found in the analyzed song
lyrics. The intertopic distance map(Figure~\ref{Figure:intertopicdistancemap}) shows
that the five topics are well separated(except for 1 and 3, that have a
significant overlap). The model has successfully captured distinct clusters of
related terms, providing valuable insight and observations:

\begin{itemize}
    \item Topic 1 appears to revolve around \textbf{existential and philosophical
      themes}, with frequent use of words like ``world``, ``life``, ``time``,
      ``die``, ``live``, and ``dream``.
    \item Topic 2 is characterized by \textbf{colloquial and explicit
      language}, exhibiting themes related to \textbf{rap and hip hop culture}.
      This topic is strongly tied to urban and street culture themes.
    \item Topic 3 revolves around themes of \textbf{love, heartbreak, emotional
      conflict and relationships}, as indicated by words like ``love``,
      ``heart``, ``hurt, ``break``.
    \item Topic 4 focuses on themes of \textbf{romantic affection, celebration,
      and festives}. Words such as ``baby``, ``tonight``, ``girl``, ``wanna",
      and ``kiss`` highlight intimate emotions and the inclusion of
      ``Christmas``, ``Santa`` and other related words strongly suggests a
      Christmas-themed music.
    \item Topic 5 appears to center around themes of \textbf{music, rhythm and
      energetic expressions}. Most prominent words identified for that topic
      are ``hey``, ``yeah``, ``rock``, ``music`` and ``beat``. They suggest a
      focus on the sound and feel of music. Inclusion of terms such as
      ``regga``(stemmed version of ``reggae``) indicates possible inclusion of
      reggae songs.
\end{itemize}

The topic modeling analysis effectively uncovered five distinct thematic
clusters within the song lyrics dataset, each representing unique aspects of
lyrical expression, ranging from love and affection, to urban culture, and
Christmas themes. The intertopic distance map
(Figure~\ref{Figure:intertopicdistancemap}) confirms that the identified topics
are well-separated, with overlaps only in related themes such as existentialism
(Topic 1) and emotional conflict (Topic 3).


\subsection{Genre Distribution Across Topics}
These results demonstrate the potential of topic modeling to extract and
analyze dominant lyrical themes in music at scale. After assigning the most
probable topic to each song the distribution of features and genres in each
topic can be checked.


\begin{center}
\begin{figure}[H]
  \centering
  \includegraphics[width=5in]{img/lda_genres_distribution.png}
  \caption{Heatmap showing percentage of genres in each topic.}
  \label{Figure:fig_eh}
\end{figure}
\end{center}

The heatmap provides a detailed representation of the percentage distribution
of genres within each topic derived from the LDA model. Key observations
include:

\begin{itemize}
  \item \textbf{Topic 1} has a balanced distribution across multiple genres,
    with slight overhead of \textit{metal}, that makes up for over 20\% of the
    songs in that topic.
  \item \textbf{Topic 2} is dominated by \textit{rap} and \textit{hip hop} genres, which
    aligns with its themes of urban culture and colloquial language.
  \item \textbf{Topic 3} focuses on love and heartbreak, a common theme in
    music across many genres. The heatmap doesn't show any dominant genre, with
    \textit{R\&B}, \textit{rock}, \textit{pop}, and \textit{country} making the
    highest contributions.
  \item \textbf{Topic 4} has a significant representation in \textit{R\&B},
    \textit{pop} and \textit{country}, suggesting its romantic and festive
    focus, including elements like celebrations and intimate moments.
  \item \textbf{Topic 5}: strongly linked to \textit{reggae} and
    \textit{rock}, reflecting the energetic and rhythmic themes of music and
    sound.
\end{itemize}

Based on the heatmap topics 1 and 3 show more diverse genre distribution,
indicating thematic universality or overlap. This observation aligns with the
conclusions drawn from the  distance
map(Figure~\ref{Figure:intertopicdistancemap}) and the analysis of their most
relevant terms((Figure~\ref{Figure:t1}, (Figure~\ref{Figure:t3})). Topics 2 and
5 display strong dominance by specific genres, showing more focused thematic
expressions.


\subsection{Empath Features Distributions Across Topics}

\begin{center}
\begin{figure}[H]
  \centering
  \includegraphics[width=5in]{img/lda_empath_features.png}
  \caption{Heatmap showing the distribution of scaled empath features across topics.}
  \label{Figure:fig_eh}
\end{figure}
\end{center}


This heatmap visualizes the scaled average Empath features by topic,
highlighting the significanace of specific lyrical themes within each topic.
Key observations include:

\begin{itemize}
  \item \textbf{Topic 1}: high values of \textit{night}, \textit{body},
    \textit{pain} and \textit{shame} suggest lyrics focusing on darker or more
    aggresive themes,  aligning well with the dominance of metal as the main
    genre.
  \item \textbf{Topic 2}: strong association with \textit{violence},
    \textit{movement}, \textit{speaking} and \textit{youth}, which are themes
    frequently occurring themes in rap and hip hop music.
  \item \textbf{Topic 3}: Strong emphasis on \textit{love}, \textit{affection},
    \textit{negative emotions}, and \textit{friends} reflects themes of
    relationships, heartbreak, and emotional conflict.
  \item \textbf{Topic 4}: High values for \textit{family}, \textit{children},
    \textit{positive emotions}, and \textit{youth} suggest themes of affection
    and relationships with loved ones, possibly relating to Christmas or other
    festive occasions.
  \item \textbf{Topic 5}: High values in \textit{music}, \textit{movement}, and
    \textit{violence} suggest rhythmic and energetic themes.
\end{itemize}

\subsection{Sentiment Across Topics}
\begin{center}
\begin{figure}[H]
  \centering
  \includegraphics[width=5in]{img/lda_sentiment.png}
  \caption{Bar chart showing average sentiment polarity by topic.}
  \label{Figure:fig_eh}
\end{figure}
\end{center}

The bar chart illustrates the average sentiment polarity for each topic. Topics
3 and 4 exhibit the highest sentiment polarity, indicating more positive
themes, while Topic 2 has the lowest polarity, reflecting more neutral or
potentially negative sentiment, which intuitively aligns with previous
findings.


\subsection{Spotify Features Across Topics}

\begin{center}
\begin{figure}[H]
  \centering
  \includegraphics[width=5in]{img/lda_spotify_features.png}
  \caption{Heatmap showing average values of scaled spotify features in each
  topic.}
  \label{Figure:fig_eh}
\end{figure}
\end{center}

The heatmap provides a visualization of the average scaled Spotify features
across topics. It shows distinct patterns for each topic, highlighting their
unique musical attributes:

\begin{itemize}
  \item \textbf{Topic 1}: High values in \textit{acousticness} and
    \textit{instrumentalness}, paired with low \textit{danceability} and
    \textit{valance} suggest calmer, more instrumental tracks.
  \item \textbf{Topic 2}: High scores in \textit{explicitness},
    \textit{speechiness}, \textit{liveness}, \textit{energy} and
    \textit{loudness} indicate tracks with s trong vocal content, high energy
    and explicit language, aligning well with the dominant genres in this
    topic.
  \item \textbf{Topic 3}: High \textit{popularity} and \textit{acousticness},
    paired with low \textit{energy} and  \textit{danceability} suggest slower,
    more emotional, acoustic tracks.
  \item \textbf{Topic 4}: High \textit{acousticness},  \textit{popularity},
    relatively high \textit{valence} and \textit{tempo} indicate s ofter,
    somewhat upbeat  tracks that include both instrumental and vocal
    components.
  \item \textbf{Topic 5}: High values in \textit{tempo}, \textit{valence}, and
    \textit{danceability}, along with moderate \textit{energy}, suggest lively,
    upbeat, and energetic and possibly danceable tracks.
\end{itemize}


\section{Temporal Trends in Music}
\label{sec:temporaltrends}

To analyze the temporal trends in musical features, during the data collection
phase the variable \textit{album release year} was stratified into 10-year
buckets representing distinct decades. This stratification allowed for a
comparative analysis of feature evolution across time periods.


\subsection{Identification of Features Affected by Release Year}

The process for identifying features influenced by release year involved two
main steps:

\begin{itemize}
  \item \textbf{Linear Regression Analysis}: Simple linear regression models
    were trained for each feature using the decade as the independent variable.
    The coefficient of determination ($R^2$) values were extracted to quantify
    how well the decade explained the variance in each feature.
  \item \textbf{ANOVA Testing}: Analysis of variance (ANOVA) was conducted to
    assess the statistical significance of differences in feature means across
    decades. The resulting p-values were used to identify features where
    temporal variation was statistically significant.
\end{itemize}
Features exhibiting high $R^2$ values and statistically significant p-values
($\alpha = 0.05$) were selected as candidates for further analysis.

24 features were identifed to  have meaningful dependency on the decade of
release, those are  their plots along with 95\% confidence intervals:

A total of 24 features were identified as having a meaningful dependency on the
decade of release. These  features demonstrated statistically significant
p-values and relatively high $R^2$ coefficients, indicating notable changes in
their values over time.


% \usepackage{color}
% \usepackage{tabularray}
\begin{table}[H]
\centering
\caption{Features identified as having significant dependency on the decade of
release.}
\begin{tblr}{
  column{even} = {r},
  column{3} = {r},
  column{5} = {r},
  hline{1,25} = {-}{0.08em},
  hline{2} = {-}{},
}
\textbf{Feature}       & \textbf{R²} & \textbf{Trend Coefficient} & \textbf{F-Statistic} & \textbf{p-Value} \\
explicit               & 0.909073    & 0.006645                   & 76.023316            & 0.000000         \\
valence                & 0.906611    & -0.003028                  & 45.487845            & 0.000000         \\
syntactic\_complexity  & 0.905622    & 0.519760                   & 3.188473             & 0.002285         \\
popularity             & 0.893476    & 0.276672                   & 61.320459            & 0.000000         \\
vader\_compound        & 0.880113    & -0.007787                  & 11.830054            & 0.000000         \\
sentiment\_polarity    & 0.853502    & -0.001479                  & 8.135521             & 0.000000         \\
loudness               & 0.833296    & 0.068308                   & 139.657491           & 0.000000         \\
rhyme\_density         & 0.824603    & -0.000094                  & 18.454532            & 0.000000         \\
danceability           & 0.814174    & 0.001072                   & 7.099057             & 0.000000         \\
speechiness            & 0.806721    & 0.000817                   & 18.774707            & 0.000000         \\
dale\_chall            & 0.775952    & 0.031306                   & 5.824913             & 0.000001         \\
flesch\_reading\_ease  & 0.771849    & -0.415172                  & 3.364072             & 0.001401         \\
gunning\_fog           & 0.758302    & 0.154463                   & 3.178465             & 0.002350         \\
lexical\_richness      & 0.565358    & -0.000709                  & 10.507191            & 0.000000         \\
type\_token\_ratio     & 0.565358    & -0.000709                  & 10.507191            & 0.000000         \\
energy                 & 0.518081    & 0.003036                   & 84.094129            & 0.000000         \\
acousticness           & 0.469885    & -0.005867                  & 221.629618           & 0.000000         \\
duration\_ms           & 0.151967    & 495.225838                 & 102.202431           & 0.000000         \\
semantic\_depth        & 0.115761    & 0.000821                   & 3.334177             & 0.001523         \\
tempo                  & 0.110478    & 0.036494                   & 2.943202             & 0.004482         \\
sentiment\_variability & 0.064628    & -0.000182                  & 5.719228             & 0.000001         \\
linguistic\_uniqueness & 0.020538    & 0.000049                   & 2.558628             & 0.012530         \\
instrumentalness       & 0.016573    & 0.000118                   & 4.748720             & 0.000025         
\end{tblr}
\end{table}


The following plots display the trends of these features across decades, along
with 95\% confidence intervals to illustrate the variability of the trends:

\begin{center}
\begin{figure}[H]
  \centering
  \includegraphics[width=5in]{img/temporal_trends_1.png}
  \caption{Plot showing temporal changes of features dependent on the decade in
  which the song was released.}
  \label{Figure:fig_eh}
\end{figure}
\end{center}

\begin{center}
\begin{figure}[H]
  \centering
  \includegraphics[width=5in]{img/temporal_trends_2.png}
  \caption{Plot showing temporal changes of features dependent on the decade in
  which the song was released.}
  \label{Figure:fig_eh}
\end{figure}
\end{center}

\begin{center}
\begin{figure}[H]
  \centering
  \includegraphics[width=5in]{img/temporal_trends_3.png}
  \caption{Plot showing temporal changes of features dependent on the decade in
  which the song was released.}
  \label{Figure:fig_eh}
\end{figure}
\end{center}

Based on those plots it can be observed that:
\begin{itemize}
  \item \textbf{There is a clear, strong upward trend in explicit content frrom
    the 1980s onward}. It reflects a cultural shift  towards more openness in
    lyrical expression and increase in popularity of genres such as rap.
  \item \textbf{There is a steady increase in popularity over the decades. Potentially
    it might reflect the increased emphasis on producing music for mass appeal.}
  \item Speechiness peaks in the 1990s and 2000s, likely reflecting the rise of
    rap and spoken-word music during this period.
  \item The observed decrease in vader compound and valence suggests a shift
    towards less positive or more emotionally nuanced lyrical content over the
    decades.
  \item \textbf{There is a substantial decrease in acousticness and a
      corresponding increase in energy over the decades, reflecting a clear
      shift in musical trends. This transition highlights the move away from
      traditional, acoustic-oriented compositions to electronically produced
      music, driven by the rise of genres like EDM and modern pop. The increase
    in energy further underscores the growing preference for high-tempo,
  dynamic, and engaging tracks.}
  \item Loudness exhibited a consistent increase over the decades, likely
    also influenced by the shift toward electronic music, which often employs
    amplified and dynamically intense production techniques.
\end{itemize}

\chapter{Conclusion}
\label{cha:conclusion}
%---------------------------------------------------------------------------

\section{Summary of Contributions}
\label{sec:summaryofcontributions}

This thesis has provided an extensive analysis of the relationships between
acoustic parameters, lyrical content, and metadata in characterizing music
tracks. By integrating well organized data collection methodology, advanced
feature engineering, explainable AI methods, machine learning models and
statistical methods, this study accomplished the following:

\begin{itemize}
  \item Constructed a diverse dataset by combining Spotify metadata and audio
    features, textual features derived from lyrics, and acoustic features
    extracted from downloaded MP3 files, representing over 3,500 songs;

  \item Employed advanced statistical methods to conduct extensive exploratory
    data analysis to identify relationships, patterns, and distinctive
    characteristics within acoustic and lyrical features across genres;

  \item Found distinctive traits of the genres included in the study,
    describing their unique acoustic and lyrical attributes;

  \item Trained machine learning models for various predictive tasks, including
    regression of popularity and  classification of explicitness, sentiment,
    and genre. and identified key features influencing these attributes. The
    models were tested against baseline models to evaluate their performance,
    and explainable AI methodologies such as SHAP were used to enhance
    interpretability and investigate feature importance and interactions;

  \item Applied Latent Dirichlet Allocation (LDA) to identify and analyze
    thematic patterns in lyrics, providing insights into the recurring themes
    that can be found in the songs included in the dataset as well as their
    distribution across genres and decades, and their distinct characteristics;

  \item With the use of statistical methods it investigated temporal trends in
    music characteristics, identifying changes reflecting the evolution of
    music styles and listener preferences across the decades;

  \item Created a reusable framework for data collection, feature engineering
    and XAI-driven analysis for music, providing foundation for further
    research in similar domains.
\end{itemize}






%---------------------------------------------------------------------------
\section{Limitations of the Study}
\label{sec:limitations}

When interpreting the results, several limitations of the  study should be
considered:

\subsubsection*{Selection Bias}

In order to acquire the data in the data collection phase the method chosen for
sampling the songs available on Spotify utilized  stratified sampling based on
genre and release decade. This method of sampling relied on Spotify's search
algorithm, and therefore introduced a potential bias influenced by factors that
this thesis did not account for. Factors like the musical preferences of the
user querying the data, the time of the year at which the data was acquired,
and many other features influencing the search algorithm's results likely
impacted the selection of songs, potentially skewing the dataset. Collection of
truly random sample for the purpose of this study would have posed significant
challenges, explained in \textit{Sampling Method} (\ref{sec:samplingmethod}).


\subsubsection*{Dimensionality}

Due to computational and memory constraints the dataset was relatively small.
A larger sample would most likely provide better quality insights and improve
the performance of the predictive models, leading to more robust findings. 


\subsubsection*{Dependency on Spotify}
Many features used in this study were derived from Spotify's metadata and
algorithms. Genre labels, popularity metrics, and features such as danceability
and loudness depend on Spotify's estimations, which may not always align with
objective or universal understanding of those properties.



%---------------------------------------------------------------------------
\section{Recommendations for Future Work}
\label{sec:recommendationsforfuturework}


\subsubsection*{Gather More Data and Increase Diversity}
Collecting more data with possibly a different sampling method and adding
more diverse features would provide a better foundation for research. It would
not only improve the quality of trained predictive models but also contribute
to more comprehensive and accurate analyses, uncovering new findings or patterns.

\subsubsection*{Reduce Dependency on Spotify}
Incorporating alternative sources into the data collection proccess would reduce
the reliance on Spotify's metadata and features generated using its models.
Possible options would be publicly available audio datasets, data from other
streaming platforms, or more extensive feature extraction from raw audio files.


\subsubsection*{Improve Feature Engineering Methodology}
The feature  engineering process could be greatly enhanced, with particular
focus on more advanced extraction of features from MP3 files. Deeper analysis
of importance of these features and more advanced feature selection would be
necessary. Improved feature engineering would give better insight into music's
structure and allow for exploration for more relationships and dependencies
with lyrical features, as well as contribute to more accurate analyses and
better performing models.

\subsubsection*{Implementation and Evaluation of More Advanced Modeling Techniques}
The study focused on model interpretability and for simplicity focused only on
using Catboost, since the goal was getting insight into model's predictive
processes and understanding relationships between features. More complex models
and architectures could likely strongly outperform the models trained in this
thesis and yield much more accurate results with the same set of features. Due
to computational and time constraints the searches with Optuna were also
relatively shallow, but the parameter grids could be much larger and numbers of
trials much higher with enough resources, allowing to find the best model
configurations.

\subsubsection*{Improve Clustering and Topic Modeling}
The clustering and topic modelling techniques could be largely extended and
provide much more information with enough effort. The topics extracted with LDA
could also serve as features in training of next models.

% \include{rozdzial8}
% \include{rozdzial9}

\cleardoublepage % Dodaj nową stronę, aby "Literatura" zaczęła się na nieparzystej stronie
\phantomsection % Dodaj punkt w spisie treści dla sekcji "Literatura"
\addcontentsline{toc}{chapter}{Bibliography} % Dodaj "Literatura" do spisu treści, bez numerowania
\printbibliography[title={Bibliography}]
\end{document}
